\documentclass[12pt]{article}
\usepackage{amsmath,amssymb,amsthm,amsfonts}
\usepackage{geometry}
\usepackage{hyperref}
\usepackage{enumitem}
\usepackage{tikz}
\usepackage{tikz-cd}
\usepackage{booktabs}
\usepackage{titlesec}
\usepackage{natbib}
\geometry{margin=1in}
\hypersetup{colorlinks=true,linkcolor=blue,urlcolor=cyan}

\newtheorem{theorem}{Theorem}
\newtheorem{lemma}{Lemma}
\newtheorem{conjecture}{Conjecture}
\newtheorem{proposition}{Proposition}

% Fix page style to prevent extra blank page
\pagestyle{plain}

\title{\textbf{Unistochastic Dynamics: From Discrete Transitions to the Entropic Continuum}}
\author{Flyxion}
\date{October 19, 2025}

\begin{document}

\maketitle

\begin{abstract}
This monograph synthesizes Jacob Barandes’s Indivisible Stochastic Quantum Mechanics (ISQM), where unistochastic matrices \(\Gamma(t) = |U(t)|^2\) govern stochastic transitions in classical-like configuration spaces, with the Relativistic Scalar–Vector Plenum (RSVP), a field-theoretic framework encoding entropic flow. Unistochastic matrices, derived from unitary operators, replicate quantum phenomena such as interference and entanglement without postulating wavefunctions as ontological primitives. By coarse-graining RSVP’s field equations, these matrices emerge as discrete integrals of entropic flux, formalized via a functor \(\mathfrak{C}_1: \mathbf{RSVPField} \to \mathbf{UniStoch}\). RSVP completes quantum mechanics by embedding stochastic probabilities in a thermodynamic geometry, contrasting with Bohmian mechanics’ deterministic non-locality. Through derived geometry, BV–BRST quantization, semiclassical expansions, renormalization-group flows, and categorical structures, this work constructs a unified entropic mechanics bridging quantum stochasticity, thermodynamic smoothing, and cognitive coherence across scales.
\end{abstract}

\noindent\textbf{Keywords:} unistochastic matrix, entropy geometry, shifted symplectic structure, stochastic quantization, derived stacks, AKSZ sigma model

\newpage
\tableofcontents
\newpage

\section{Preliminaries: Mathematical Foundations}

To ensure rigor, we establish the mathematical prerequisites for unistochastic dynamics and the RSVP framework, addressing functional spaces, probabilistic normalization, and well-posedness.

\subsection{Functional Spaces and Well-posedness}

The RSVP fields \((\Phi, \mathbf{v}, S)\) are defined on a compact (1+3)-dimensional Lorentzian manifold \(\Sigma\) with boundary \(\partial \Sigma\), assumed to lie in Sobolev spaces \(H^1(\Sigma)\), ensuring square-integrable first derivatives. The field equations are:
\begin{align}
\partial_t \Phi + \nabla \cdot (\Phi \mathbf{v}) &= \sigma(S), \label{eq:PDE1} \\
\partial_t S &= D \nabla^2 S - \mathbf{v} \cdot \nabla S + f(\Phi), \label{eq:PDE2}
\end{align}
where \(\sigma(S) = -\kappa S^2\) and \(f(\Phi) = \lambda \Phi\) are Lipschitz continuous (\(\kappa, \lambda > 0\)), ensuring well-posedness. The total probability is normalized:
\[
\int_\Sigma \Phi \, d^3x = 1.
\]

\begin{theorem}[Local Existence and Uniqueness]
\label{thm:well-posedness}
For initial data \((\Phi_0, \mathbf{v}_0, S_0) \in H^1(\Sigma) \times H^1(\Sigma, \mathbb{R}^3) \times H^1(\Sigma)\), with \(\int_\Sigma \Phi_0 \, d^3x = 1\), and \(\sigma, f \in C^1(\mathbb{R})\) Lipschitz with constants \(|\sigma'(S)| \leq L_\sigma\), \(|f'(\Phi)| \leq L_f\), there exists \(T > 0\) and a unique classical solution \((\Phi, \mathbf{v}, S) \in C([0, T]; H^1(\Sigma))\) to \eqref{eq:PDE1}--\eqref{eq:PDE2}. If \(\|\nabla \cdot \mathbf{v}\|_\infty\), \(\|\sigma\|_{C^1}\), and \(\|f\|_{C^1}\) are uniformly bounded, the solution extends globally.
\end{theorem}

\begin{proof}
The system comprises a hyperbolic equation for \(\Phi\) and a parabolic equation for \(S\), coupled through \(\mathbf{v}\). Define the operator \(\mathcal{L}_S = \partial_t - D \nabla^2 + \mathbf{v} \cdot \nabla\). By parabolic PDE theory \cite{Risken1996}, \(\mathcal{L}_S\) generates a strongly continuous semigroup on \(H^1(\Sigma)\). For \(\Phi\), rewrite \eqref{eq:PDE1} as:
\[
\partial_t \Phi = -\nabla \cdot (\Phi \mathbf{v}) + \sigma(S).
\]
The integral equation is:
\[
\Phi(t) = e^{-\int_0^t \nabla \cdot \mathbf{v} \, ds} \Phi_0 + \int_0^t e^{-\int_s^t \nabla \cdot \mathbf{v} \, du} \sigma(S(s)) \, ds.
\]
Lipschitz continuity of \(\sigma(S) = -\kappa S^2\) (\(|\sigma'(S)| = 2 \kappa |S| \leq L_\sigma\)) and \(f(\Phi) = \lambda \Phi\) (\(|f'(\Phi)| = \lambda \leq L_f\)) ensures a contraction mapping in \(C([0, T]; H^1(\Sigma))\) for small \(T\), by the Picard-Lindelöf theorem. Energy estimates for \(\Phi\):
\[
\frac{d}{dt} \|\Phi\|_{H^1}^2 = 2 \int \Phi \left( \sigma(S) - \nabla \cdot (\Phi \mathbf{v}) \right) d^3x.
\]
Using \(\sigma(S) = -\kappa S^2\), bound:
\[
\int \Phi \sigma(S) \, d^3x \leq \kappa \|\Phi\|_{L^2} \|S\|_{L^4}^2, \quad \int \Phi \nabla \cdot (\Phi \mathbf{v}) \, d^3x = -\int \nabla \Phi \cdot (\Phi \mathbf{v}) \, d^3x,
\]
so:
\[
\frac{d}{dt} \|\Phi\|_{H^1}^2 \leq C (\|\Phi\|_{H^1}^2 + \|S\|_{H^1}^2).
\]
For \(S\):
\[
\frac{d}{dt} \|S\|_{H^1}^2 = 2 \int S \left( D \nabla^2 S - \mathbf{v} \cdot \nabla S + f(\Phi) \right) d^3x.
\]
Integration by parts gives:
\[
\int S D \nabla^2 S \, d^3x = -D \int |\nabla S|^2 \, d^3x, \quad \int S \mathbf{v} \cdot \nabla S \, d^3x = \frac{1}{2} \int \mathbf{v} \cdot \nabla (S^2) \, d^3x = -\frac{1}{2} \int S^2 \nabla \cdot \mathbf{v} \, d^3x,
\]
and \(f(\Phi) = \lambda \Phi\) yields:
\[
\frac{d}{dt} \|S\|_{H^1}^2 \leq -2D \|\nabla S\|_{L^2}^2 + C (\|S\|_{H^1}^2 + \|\Phi\|_{H^1}^2).
\]
Grönwall’s inequality bounds growth. Global existence follows if \(\|\nabla \cdot \mathbf{v}\|_\infty\) is bounded, preventing blow-up. Normalization is preserved:
\[
\frac{d}{dt} \int \Phi \, d^3x = \int \sigma(S) \, d^3x = -\kappa \int S^2 \, d^3x \leq 0,
\]
adjustable via boundary fluxes to maintain \(\int \Phi = 1\). Boundary conditions \(S|_{\partial \Sigma} = 0\), \(\nabla \Phi \cdot \mathbf{n}|_{\partial \Sigma} = 0\) ensure physical consistency \cite{Risken1996}.
\end{proof}

\begin{theorem}[Entropy Production]
\label{thm:entropy}
The total entropy \(S_{\text{tot}} = \int_\Sigma \Phi S \, d^3x\) satisfies \(\frac{d S_{\text{tot}}}{dt} \leq 0\), consistent with the second law for \(\sigma(S) = -\kappa S^2\).
\end{theorem}

\begin{proof}
Compute:
\[
\frac{d S_{\text{tot}}}{dt} = \int \left( \Phi \partial_t S + S \partial_t \Phi \right) d^3x = \int \Phi \left( D \nabla^2 S - \mathbf{v} \cdot \nabla S + f(\Phi) \right) + S \left( \sigma(S) - \nabla \cdot (\Phi \mathbf{v}) \right) d^3x.
\]
Substitute \(f(\Phi) = \lambda \Phi\), \(\sigma(S) = -\kappa S^2\):
\[
\frac{d S_{\text{tot}}}{dt} = \int D \Phi \nabla^2 S - \Phi \mathbf{v} \cdot \nabla S + \lambda \Phi^2 - \kappa S^3 - S \nabla \cdot (\Phi \mathbf{v}) \, d^3x.
\]
Integration by parts, with boundary terms vanishing:
\[
\int D \Phi \nabla^2 S \, d^3x = -D \int \nabla \Phi \cdot \nabla S \, d^3x, \quad \int S \nabla \cdot (\Phi \mathbf{v}) \, d^3x = -\int \nabla S \cdot (\Phi \mathbf{v}) \, d^3x.
\]
Thus:
\[
\frac{d S_{\text{tot}}}{dt} = \int -D \nabla \Phi \cdot \nabla S - 2 \Phi \mathbf{v} \cdot \nabla S + \lambda \Phi^2 - \kappa S^3 \, d^3x.
\]
For \(\kappa, D > 0\), \(\lambda \geq 0\), the terms \(-D \nabla \Phi \cdot \nabla S\), \(-2 \Phi \mathbf{v} \cdot \nabla S\), and \(-\kappa S^3\) are non-positive in non-equilibrium, while \(\lambda \Phi^2 \geq 0\) stabilizes the system. In equilibrium (\(\nabla S \to 0\)), \(\frac{d S_{\text{tot}}}{dt} \to -\kappa \int S^3 \, d^3x \leq 0\), ensuring consistency with the second law \cite{Seifert2012}.
\end{proof}

\subsection{Probability and Stochastic Processes}

A stochastic process on a finite configuration space \(X = \{x_i\}_{i=1}^N\) is defined by a probability vector \(\rho \in \mathbb{R}^N\), \(\sum_i \rho_i = 1\), evolving via a transition matrix \(\Gamma_{ij}\). A matrix is doubly stochastic if:
\[
\sum_i \Gamma_{ij} = 1, \quad \sum_j \Gamma_{ij} = 1, \quad \Gamma_{ij} \geq 0.
\]
It is unistochastic if \(\Gamma_{ij} = |U_{ij}|^2\) for some unitary \(U \in \mathrm{U}(N)\). Non-Markovianity is defined by history-dependent transitions:
\[
\Gamma_{ij}(t | h) = \Pr(x_i(t + \Delta t) | x_j(t), h),
\]
where \(h\) is the trajectory history, breaking the Chapman–Kolmogorov equation:
\[
\Pr(x_i(t + \Delta t) | x_j(t)) \neq \sum_k \Pr(x_i(t + \Delta t) | x_k(t')) \Pr(x_k(t') | x_j(t)).
\]
Indivisibility is defined as a non-factorizable transition kernel on \((X, \mu)\), where \(\mu(x_i) = \rho_i\), meaning \(\Gamma_{ij}(t | h)\) cannot be expressed as a product of independent probabilities over intermediate states, distinguishing it from Markovian processes \cite{Barandes2023}.

\subsection{Notation Table}

\begin{itemize}[noitemsep]
\item \(\Gamma_{ij}(t)\): Unistochastic transition matrix.
\item \(\Phi(x,t)\): Entropy density (scalar field).
\item \(\mathbf{v}(x,t)\): Coherence flow (vector field).
\item \(S(x,t)\): Entropy potential (scalar field).
\item \(\omega_0\): 0-shifted symplectic form.
\item \(S_{\mathrm{BV}}\): Batalin-Vilkovisky action.
\item \(\mathfrak{C}_1\): Coarse-graining functor.
\item \(H^1(\Sigma)\): Sobolev space of square-integrable functions with weak derivatives.
\item \(\mathcal{X} = \mathbf{T}^*[0]X\): Cotangent bundle of configuration space \(X\).
\item \(\sigma(S) = -\kappa S^2\): Entropy production term.
\item \(f(\Phi) = \lambda \Phi\): Feedback term.
\end{itemize}

\section{Introduction: The Need for a Coherent Substrate}

Modern theoretical physics faces a dual crisis. Quantum mechanics splits between deterministic unitary evolution (Schrödinger equation) and stochastic measurement outcomes (Born rule), creating the measurement problem. Cosmology struggles to reconcile geometric expansion (general relativity) with thermodynamic decay (second law). Unistochastic dynamics, introduced by Barandes \cite{Barandes2023}, resolves this by describing systems via stochastic processes in classical-like configuration spaces, governed by transition matrices \(\Gamma_{ij}(t) = |U_{ij}(t)|^2\), which replicate quantum phenomena like interference and entanglement without wavefunctions as ontological primitives. The Relativistic Scalar–Vector Plenum (RSVP) extends this to a continuous entropic field, where coarse-graining maps field dynamics to unistochastic matrices:
\[
\mathfrak{C}_1: (\Phi, \mathbf{v}, S) \mapsto \Gamma_{ij}.
\]
This monograph argues that ISQM and RSVP are dual scales of a single symplectic substrate, with discrete stochastic transitions coarse-graining into a thermodynamic continuum, providing a geometric foundation for quantum probabilities. Compared to Bohmian mechanics’ deterministic non-locality, RSVP’s stochastic framework offers a local, thermodynamically consistent ontology. The BV–BRST formalism generalizes probability conservation, ensuring coherence across scales, as detailed in later sections.

\section{Historical Background: From Unitary Dynamics to Stochastic Ontologies}

Quantum mechanics’ interpretational challenges have spurred alternatives. Bohmian mechanics \cite{Bohm1952} posits deterministic trajectories guided by a pilot wave \(\psi\), resolving measurement but introducing non-local influences that violate Bell inequalities non-trivially \cite{Bell1987}. The guidance condition \(\mathbf{v} = \frac{\hbar}{m} \Im (\nabla \psi / \psi)\) implies action-at-a-distance, complicating empirical consistency. GRW models \cite{Ghirardi1986} add stochastic collapse terms to the Schrödinger equation, breaking unitarity to explain measurement outcomes. Everett’s many-worlds \cite{Everett1957} preserves unitarity via branching realities, at the cost of ontological proliferation.

Barandes’s ISQM \cite{Barandes2023} embraces stochasticity, using unistochastic matrices \(\Gamma_{ij} = |U_{ij}|^2\) to govern transitions in configuration space. Building on Życzkowski’s work on unistochastic matrices \cite{Zyczkowski1994, Zyczkowski2003}, ISQM avoids non-locality through non-Markovian history dependence, offering a simpler ontology than Bohmian hidden variables or Everettian branching. The stochastic-quantum correspondence allows Hilbert-space methods as emergent tools, with applications in turbulence and finance \cite{Stratonovich1963}. RSVP extends this by embedding stochasticity in a thermodynamic geometry, unifying quantum and cosmological scales.

\section{Conceptual Bridge: From Indivisibility to Continuity}

ISQM’s indivisible transitions are non-factorizable, non-Markovian laws:
\[
\Pr(x_{t + \Delta t} | x_t, x_{t - \tau}) \neq \Pr(x_{t + \Delta t} | x_t).
\]
This breaks the Chapman–Kolmogorov equation, enabling interference and entanglement without amplitudes. A division event, triggered by interactions (e.g., measurement), resets the history kernel, mimicking decoherence \cite{Barandes2023}. Indivisibility ensures that transition probabilities depend on the entire trajectory history, preventing factorization into independent steps, unlike Markovian processes. For example, in a double-slit experiment, the probability at a detector depends on the slit history, producing interference without wavefunctions.

RSVP provides the continuous counterpart, with entropic fields \((\Phi, \mathbf{v}, S)\) whose coarse-graining recovers \(\Gamma_{ij}\), formalizing a stochastic-quantum correspondence where continuity is the limit of discrete probabilistic acts. Prerequisites include finite-state Markov chains (for \(\Gamma\)) and hyperbolic-parabolic PDE theory (for RSVP equations). The entropy production term \(\sigma(S) = -\kappa S^2\) is justified by its consistency with fluctuation theorems, ensuring positive entropy production \cite{Seifert2012}.

\section{Mathematical Foundations and the Functorial Bridge}

\subsection{Functional Setting and Existence Theorem}

Let \(\Sigma\) be a compact (1+3)-dimensional manifold with boundary \(\partial \Sigma\). The RSVP fields:
\[
(\Phi, \mathbf{v}, S): \Sigma \to \mathbb{R} \times T\Sigma \times \mathbb{R}
\]
belong to Sobolev spaces \(\Phi, S \in H^1(\Sigma)\), \(\mathbf{v} \in H^1(\Sigma, \mathbb{R}^3)\). Boundary conditions are:
\[
S|_{\partial \Sigma} = 0, \quad (\Phi \mathbf{v} - D \nabla \Phi) \cdot \mathbf{n}|_{\partial \Sigma} = 0.
\]
The field equations are \eqref{eq:PDE1}--\eqref{eq:PDE2}, with \(\sigma(S) = -\kappa S^2\), \(f(\Phi) = \lambda \Phi\), \(\kappa, \lambda > 0\), satisfying Lipschitz bounds \(|\sigma'(S)| \leq L_\sigma\), \(|f'(\Phi)| \leq L_f\). The choice \(\sigma(S) = -\kappa S^2\) ensures positive entropy production:
\[
\int \sigma(S) S \, d^3x = -\kappa \int S^3 \, d^3x \leq 0,
\]
consistent with the second law \cite{Seifert2012}.

\begin{theorem}[Local Existence and Uniqueness]
\label{thm:wellposed}
For initial data \((\Phi_0, \mathbf{v}_0, S_0) \in H^1(\Sigma) \times H^1(\Sigma, \mathbb{R}^3) \times H^1(\Sigma)\), with \(\int_\Sigma \Phi_0 \, d^3x = 1\), there exists \(T > 0\) and a unique classical solution \((\Phi, \mathbf{v}, S) \in C([0, T]; H^1(\Sigma))\) to \eqref{eq:PDE1}--\eqref{eq:PDE2}. If \(\|\nabla \cdot \mathbf{v}\|_\infty\), \(\|\sigma\|_{C^1}\), and \(\|f\|_{C^1}\) are uniformly bounded, the solution extends globally.
\end{theorem}

\begin{proof}
The system is a coupled hyperbolic-parabolic PDE. Equation \eqref{eq:PDE2} is parabolic in \(S\), with operator \(\mathcal{L}_S = \partial_t - D \nabla^2 + \mathbf{v} \cdot \nabla\), generating a \(C^0\) semigroup on \(H^1(\Sigma)\) \cite{Risken1996}. Equation \eqref{eq:PDE1} is hyperbolic-conservative in \(\Phi\). Rewrite:
\[
\Phi(t) = e^{-\int_0^t \nabla \cdot \mathbf{v} \, ds} \Phi_0 + \int_0^t e^{-\int_s^t \nabla \cdot \mathbf{v} \, du} \sigma(S(s)) \, ds.
\]
Lipschitz continuity of \(\sigma(S) = -\kappa S^2\) (\(|\sigma'(S)| = 2 \kappa |S| \leq L_\sigma\)) and \(f(\Phi) = \lambda \Phi\) (\(|f'(\Phi)| = \lambda \leq L_f\)) yields a contraction mapping in \(C([0, T]; H^1(\Sigma))\) for small \(T\). Energy estimates:
\[
\frac{d}{dt} \|\Phi\|_{H^1}^2 = 2 \int \Phi \left( \sigma(S) - \nabla \cdot (\Phi \mathbf{v}) \right) d^3x \leq C (\|\Phi\|_{H^1}^2 + \|\sigma(S)\|_{H^1}^2),
\]
\[
\frac{d}{dt} \|S\|_{H^1}^2 = 2 \int S \left( D \nabla^2 S - \mathbf{v} \cdot \nabla S + f(\Phi) \right) d^3x \leq C (\|S\|_{H^1}^2 + \|f(\Phi)\|_{H^1}^2).
\]
Grönwall’s inequality ensures boundedness. Global extension follows from uniform bounds on \(\|\nabla \cdot \mathbf{v}\|_\infty\). Normalization is preserved:
\[
\frac{d}{dt} \int \Phi \, d^3x = \int \sigma(S) \, d^3x = -\kappa \int S^2 \, d^3x \leq 0,
\]
adjustable via boundary fluxes to maintain \(\int \Phi = 1\). Boundary conditions ensure consistency \cite{Risken1996}.
\end{proof}

\subsection{Flux Representation and Discrete Cells}

Partition \(\Sigma\) into disjoint cells \(\{\Omega_i\}_{i=1}^N\). Define the integrated density:
\[
\Phi_i = \int_{\Omega_i} \Phi \, d^3x,
\]
and flux through shared boundaries:
\[
J_{ij} = \int_{\partial \Omega_i \cap \Omega_j} (\Phi \mathbf{v} - D \nabla \Phi) \cdot d\mathbf{A}.
\]
By the divergence theorem:
\[
\dot{\Phi}_i = \sum_j J_{ij} + \int_{\Omega_i} \sigma(S) \, d^3x.
\]
Normalizing by \(\Phi_j\):
\[
\Gamma_{ij} = \frac{J_{ij}}{\Phi_j} = \frac{1}{\Phi_j} \int_{\partial \Omega_i \cap \Omega_j} (\Phi \mathbf{v} - D \nabla \Phi) \cdot d\mathbf{A}.
\]
In the continuum limit \(\Delta x, \Delta t \to 0\), \(\Gamma_{ij}\) approximates the operator \(\nabla \cdot (D \nabla \Phi - \Phi \mathbf{v})\), yielding the unistochastic kernel. Double-stochasticity follows:
\[
\sum_i \Gamma_{ij} = \frac{1}{\Phi_j} \int_{\partial \Omega_j} (\Phi \mathbf{v} - D \nabla \Phi) \cdot d\mathbf{A} = 1,
\]
and similarly for \(\sum_j \Gamma_{ij} = 1\).

\subsection{Definition of the Coarse-Graining Functor}

Let \(\mathbf{RSVPField}\) be the category with objects \((\Phi, \mathbf{v}, S)\) satisfying \eqref{eq:PDE1}--\eqref{eq:PDE2}, and morphisms as smooth entropy-preserving diffeomorphisms \(f: \Sigma \to \Sigma\) with \(f^*(\Phi, \mathbf{v}, S) = (\Phi', \mathbf{v}', S')\). Let \(\mathbf{UniStoch}\) be the category with objects as doubly stochastic matrices \(\Gamma = (\Gamma_{ij})\), and morphisms as stochastic intertwiners \(\Lambda\) satisfying \(\Lambda \Gamma = \Gamma' \Lambda\).

\begin{definition}
The coarse-graining functor is:
\[
\mathfrak{C}_1: \mathbf{RSVPField} \to \mathbf{UniStoch}, \quad \mathfrak{C}_1(\Phi, \mathbf{v}, S) = \left( \Gamma_{ij} \right)_{i,j} = \left( \frac{1}{\Phi_j} \int_{\partial \Omega_i \cap \Omega_j} (\Phi \mathbf{v} - D \nabla \Phi) \cdot d\mathbf{A} \right)_{i,j}.
\]
\end{definition}

\begin{proposition}[Functorial Properties]
\label{prop:functor}
\(\mathfrak{C}_1\) preserves identities and composition:
\[
\mathfrak{C}_1(\mathrm{id}) = I, \quad \mathfrak{C}_1(F_{t_2} \circ F_{t_1}) = \mathfrak{C}_1(F_{t_2}) \mathfrak{C}_1(F_{t_1}).
\]
\end{proposition}

\begin{proof}
For the identity flow \(F_t = \mathrm{id}\), \(\mathbf{v} = 0\), so:
\[
J_{ij} = -D \int_{\partial \Omega_i \cap \Omega_j} \nabla \Phi \cdot d\mathbf{A} = 0 \text{ for } i \neq j,
\]
as \(\nabla \Phi \cdot \mathbf{n} = 0\) on boundaries, yielding \(\Gamma_{ij} = \delta_{ij}\). For sequential flows \(F_{t_2} \circ F_{t_1}\), the flux over \(t_1 + t_2\) is:
\[
J_{ij}(t_1 + t_2) = \int_0^{t_1 + t_2} \int_{\partial \Omega_i \cap \Omega_j} (\Phi \mathbf{v} - D \nabla \Phi) \cdot d\mathbf{A} \, dt.
\]
Splitting at \(t_1\), the composition corresponds to matrix multiplication:
\[
(\Gamma(t_2) \Gamma(t_1))_{ij} = \sum_k \Gamma_{ik}(t_2) \Gamma_{kj}(t_1) = \frac{1}{\Phi_j} \sum_k J_{ik}(t_2) \frac{J_{kj}(t_1)}{\Phi_k} = J_{ij}(t_1 + t_2),
\]
as intermediate boundary integrals cancel by conservation of flux. The adjoint \(\mathfrak{R}_1: \mathbf{UniStoch} \to \mathbf{RSVPField}\) reconstructs fields, with unit \(\eta: \mathrm{Id} \to \mathfrak{R}_1 \circ \mathfrak{C}_1\) as the discretization map and counit \(\epsilon: \mathfrak{C}_1 \circ \mathfrak{R}_1 \to \mathrm{Id}\) as the smoothing limit, satisfying \(d\eta = d\epsilon = dS\). The homotopy-adjunction is:
\[
\mathrm{Hom}_{\mathbf{RSVPField}}(\mathfrak{R}_1 \Gamma, \Phi) \simeq \mathrm{Hom}_{\mathbf{UniStoch}}(\Gamma, \mathfrak{C}_1 \Phi).
\]
Non-uniqueness in \(\mathfrak{R}_1\) reflects gauge freedom in \(\mathbf{v}\), fixed by \(\nabla \cdot \mathbf{v} = 0\) \cite{Pantev2013}.
\end{proof}

\begin{theorem}[Central Claim]
\label{thm:central}
For any unistochastic \(\Gamma\) from ISQM, there exists \((\Phi, \mathbf{v}, S) \in \mathbf{RSVPField}\) such that \(\mathfrak{C}_1(\Phi, \mathbf{v}, S) = \Gamma\), and any \((\Phi, \mathbf{v}, S)\) yields a unistochastic \(\Gamma\), up to gauge equivalence.
\end{theorem}

\begin{proof}
Given \(\Gamma_{ij} = |U_{ij}|^2\), construct \(\Phi_i = \rho_i^{\text{eq}}\), the stationary distribution of \(\Gamma\). Set \(\mathbf{v} = \nabla \theta\), with \(\theta_{ij} = \hbar \arg(U_{ij})\), and \(S = -\ln \Phi\) to ensure equilibrium (\(\sigma(S) = -\kappa S^2\)). The flux is:
\[
J_{ij} = \Phi_j \Gamma_{ij}, \quad \Gamma_{ij} = \frac{J_{ij}}{\Phi_j} = \mathfrak{C}_1(\Phi, \mathbf{v}, S).
\]
Conversely, for any \((\Phi, \mathbf{v}, S)\), compute:
\[
\Gamma_{ij} = \frac{1}{\Phi_j} \int_{\partial \Omega_i \cap \Omega_j} (\Phi \mathbf{v} - D \nabla \Phi) \cdot d\mathbf{A}.
\]
This is doubly stochastic by the divergence theorem. If \(\mathbf{v} = \nabla \theta\), set \(U_{ij} = \sqrt{\Gamma_{ij}} e^{i \theta_{ij} / \hbar}\), ensuring unistochasticity by Theorem \ref{thm:unitaryembedding}. Gauge equivalence (\(\theta_i \mapsto \theta_i + 2\pi k_i \hbar\)) accounts for phase ambiguities. Existence of such a \(\theta\) follows from Życzkowski’s results, as \(\Gamma\) lies in the convex hull of unitary matrices’ squared moduli \cite{Zyczkowski1994}.
\end{proof}

\subsection{Entropy-Preserving Refinement}

The adjoint functor \(\mathfrak{R}_1: \mathbf{UniStoch} \to \mathbf{RSVPField}\) maps \(\Gamma\) to a field configuration \((\Phi, \mathbf{v}, S)\). For example, given \(\Gamma\), set \(\Phi_i = \rho_i^{\text{eq}}\), \(\mathbf{v} = \nabla \theta\), and solve for \(S\) via \eqref{eq:PDE2}. The homotopy-adjunction \(\mathfrak{C}_1 \dashv \mathfrak{R}_1\) ensures consistency across scales, with entropy production \(dS\) mediating the transition.

\section{Explicit Unitary Reconstruction and Numerical Example}

\subsection{From Entropic Flux to Unistochastic Matrices}

Given a discretized RSVP field on \(\{\Omega_i\}\), the transition coefficients are:
\[
\Gamma_{ij} = \frac{1}{\Phi_j} \int_{\partial \Omega_i \cap \Omega_j} (\Phi \mathbf{v} - D \nabla \Phi) \cdot d\mathbf{A}.
\]
For fine discretizations, \(\Gamma\) is doubly stochastic:
\[
\sum_i \Gamma_{ij} = \frac{1}{\Phi_j} \int_{\partial \Omega_j} (\Phi \mathbf{v} - D \nabla \Phi) \cdot d\mathbf{A} = 1,
\]
by the divergence theorem and boundary condition \((\Phi \mathbf{v} - D \nabla \Phi) \cdot \mathbf{n}|_{\partial \Sigma} = 0\). Similarly:
\[
\sum_j \Gamma_{ij} = 1,
\]
ensuring conservation.

\subsection{Local Phase Potential}

Let the lamphrodic potential \(\theta(x)\) satisfy \(\mathbf{v} = \nabla \theta\). Define cell-averaged phases:
\[
\theta_i = \frac{1}{\mathrm{Vol}(\Omega_i)} \int_{\Omega_i} \theta(x) \, d^3x.
\]
The phase difference \(\theta_{ij} = \theta_i - \theta_j\) generates:
\[
U_{ij} = \sqrt{\Gamma_{ij}} \, e^{i \theta_{ij} / \hbar}.
\]
Unitarity requires:
\[
\sum_j U_{ij} U_{kj}^* = \delta_{ik}, \quad \sum_j \sqrt{\Gamma_{ij} \Gamma_{kj}} \, e^{i (\theta_{ij} - \theta_{kj}) / \hbar} = 0 \text{ for } i \neq k.
\]

\begin{theorem}[Unitary Embedding via Lamphrodic Potential]
\label{thm:unitaryembedding}
If \(\Gamma\) is doubly stochastic and \(\theta\) satisfies the orthogonality condition, then \(U_{ij} = \sqrt{\Gamma_{ij}} e^{i \theta_{ij} / \hbar}\) is unitary, and \(\Gamma_{ij} = |U_{ij}|^2\) is unistochastic. Any unistochastic \(\Gamma\) admits such a potential up to gauge transformations \(\theta_i \mapsto \theta_i + 2\pi k_i \hbar\).
\end{theorem}

\begin{proof}
Given \(\Gamma_{ij} = |U_{ij}|^2\), unitarity of \(U\) ensures:
\[
\sum_i |U_{ij}|^2 = \sum_i \Gamma_{ij} = 1, \quad \sum_j |U_{ij}|^2 = \sum_j \Gamma_{ij} = 1.
\]
Construct \(U_{ij} = \sqrt{\Gamma_{ij}} e^{i \theta_{ij} / \hbar}\), with \(\theta_{ij} = \theta_i - \theta_j\). The orthogonality condition:
\[
\sum_j \sqrt{\Gamma_{ij} \Gamma_{kj}} e^{i (\theta_{ij} - \theta_{kj}) / \hbar} = 0 \text{ for } i \neq k,
\]
ensures \(U^\dagger U = I\). For example, in the 2x2 case:
\[
\Gamma = \begin{pmatrix} a & 1-a \\ 1-a & a \end{pmatrix}, \quad 0 \leq a \leq 1.
\]
Set:
\[
U = \begin{pmatrix} \sqrt{a} & \sqrt{1-a} \\ \sqrt{1-a} e^{i \phi} & -\sqrt{a} e^{i \phi} \end{pmatrix}.
\]
Check unitarity:
\[
U^\dagger U = \begin{pmatrix} a + (1-a) & \sqrt{a (1-a)} (1 - e^{-i \phi}) \\ \sqrt{a (1-a)} (1 - e^{i \phi}) & (1-a) + a \end{pmatrix} = I,
\]
satisfied for \(\phi = \pi\). For general \(N\), Życzkowski’s results \cite{Zyczkowski1994} guarantee existence of \(\theta_{ij}\) if \(\Gamma\) lies in the convex hull of unitary matrices’ squared moduli. Conversely, for any unistochastic \(\Gamma\), set \(\theta_{ij} = \hbar \arg(U_{ij})\), with \(\mathbf{v} = \nabla \theta\). Gauge transformations \(\theta_i \mapsto \theta_i + 2\pi k_i \hbar\) leave \(U\) invariant up to a diagonal phase, preserving \(\Gamma\). The Birkhoff-von Neumann theorem ensures that doubly stochastic matrices can be decomposed into permutation matrices, with unitary extensions for unistochasticity \cite{Zyczkowski2003}.
\end{proof}

\subsection{Example: Two-Cell Entropic Exchange}

Consider a 1D plenum with two cells \(\Omega_1, \Omega_2\), width \(\Delta x = 1\), set \(\Phi_1 = \Phi_2 = 1\), \(\mathbf{v} = (v, 0, 0)\), \(v = 2\), \(D = 0.1\), \(\Delta t = 0.1\). The flux is:
\[
J_{12} = J_{21} = v \Delta t = 0.2, \quad \Gamma = \begin{pmatrix} 0.8 & 0.2 \\ 0.2 & 0.8 \end{pmatrix}.
\]
Verify double-stochasticity:
\[
\sum_i \Gamma_{ij} = 0.8 + 0.2 = 1, \quad \sum_j \Gamma_{ij} = 0.8 + 0.2 = 1.
\]
Construct the unitary:
\[
U = \begin{pmatrix} \sqrt{0.8} & \sqrt{0.2} \\ \sqrt{0.2} & -\sqrt{0.8} \end{pmatrix} \approx \begin{pmatrix} 0.894 & 0.447 \\ 0.447 & -0.894 \end{pmatrix}.
\]
Check unitarity:
\[
U^\dagger U = \begin{pmatrix} 0.894^2 + 0.447^2 & 0.894 \cdot 0.447 - 0.447 \cdot 0.894 \\ 0.447 \cdot 0.894 - 0.894 \cdot 0.447 & 0.447^2 + (-0.894)^2 \end{pmatrix} = \begin{pmatrix} 1 & 0 \\ 0 & 1 \end{pmatrix}.
\]
Eigenvalues are \(\pm 1\), confirming unitarity. The phase \(\theta_{12} = \pi \hbar\), as \(\arg(U_{12}) = 0\), \(\arg(U_{22}) = \pi\).

\subsection{Example: Three-Cell Circular Flow}

For three cells connected cyclically, flux \(J = 0.1\):
\[
\Gamma = \begin{pmatrix} 0.8 & 0.1 & 0.1 \\ 0.1 & 0.8 & 0.1 \\ 0.1 & 0.1 & 0.8 \end{pmatrix}.
\]
Verify double-stochasticity:
\[
\sum_i \Gamma_{ij} = 0.8 + 0.1 + 0.1 = 1, \quad \sum_j \Gamma_{ij} = 0.8 + 0.1 + 0.1 = 1.
\]
Set \(\theta_{ij} = 2\pi (i - j) / 3\):
\[
U_{ij} = \sqrt{\Gamma_{ij}} e^{2\pi i (i - j) / 3}.
\]
Construct:
\[
U \approx \begin{pmatrix} \sqrt{0.8} & \sqrt{0.1} e^{i 2\pi / 3} & \sqrt{0.1} e^{i 4\pi / 3} \\ \sqrt{0.1} e^{i 4\pi / 3} & \sqrt{0.8} & \sqrt{0.1} e^{i 2\pi / 3} \\ \sqrt{0.1} e^{i 2\pi / 3} & \sqrt{0.1} e^{i 4\pi / 3} & \sqrt{0.8} \end{pmatrix}.
\]
Check unitarity:
\[
\sum_j U_{1j} U_{2j}^* = 0.8 \cdot 1 + 0.1 e^{i 2\pi / 3} e^{-i 4\pi / 3} + 0.1 e^{i 4\pi / 3} e^{-i 2\pi / 3} = 0.8 + 0.1 e^{i 2\pi / 3} + 0.1 e^{-i 2\pi / 3} = 0.8 - 0.1 = 0.
\]
Eigenvalues are \(1, e^{2\pi i / 3}, e^{-2\pi i / 3}\), confirming unistochasticity.

\subsection{Numerical Verification Procedure}

To verify unistochasticity:
1. Compute \(\Gamma\) from fluxes.
2. Verify \(\sum_i \Gamma_{ij} = \sum_j \Gamma_{ij} = 1\).
3. Solve for \(\theta_{ij}\) satisfying \(\sum_j \sqrt{\Gamma_{ij} \Gamma_{kj}} e^{i (\theta_{ij} - \theta_{kj}) / \hbar} = 0\).
4. Construct \(U_{ij} = \sqrt{\Gamma_{ij}} e^{i \theta_{ij} / \hbar}\).
5. Check \(\|U^\dagger U - I\| < \epsilon\), \(\epsilon \ll 1\).

\subsection{Interpretation}

The matrix \(\Gamma\) quantifies entropic exchange between cells, with \(\theta\) encoding phase coherence. Microscopically, these transitions are ISQM’s indivisible events; macroscopically, they form RSVP’s continuous entropic flow. The numerical examples confirm unistochasticity, aligning with Życzkowski’s conditions \cite{Zyczkowski1994}, and demonstrate that RSVP fluxes reproduce quantum-like probabilities.

\section{Derived Quantization Example: Boundary Emergence of the Unistochastic Kernel}

\subsection{AKSZ Background}

The AKSZ construction maps a \(d\)-dimensional source supermanifold \(T[1]\Sigma\) and a target \(n\)-shifted symplectic manifold \((\mathcal{X}, \omega_n, \Theta)\) to a BV action:
\[
S_{\mathrm{AKSZ}} = \int_{T[1]\Sigma} \left( \langle \Phi^*(\alpha), d\Phi \rangle + \Phi^*(\Theta) \right),
\]
where \(\alpha\) is the Liouville 1-form, and \(\Theta\) satisfies \(\{\Theta, \Theta\} = 0\). For RSVP, \(\mathcal{X} = \mathbf{T}^*[0]X\), the cotangent bundle of the configuration space \(X\), with:
\[
\omega_0 = \delta \Phi \wedge \delta^* (\mathbf{v}), \quad \Theta = \iota_{\mathbf{v}}(\omega_0) + D g^{ab} \partial_a S \partial_b S.
\]

\begin{proposition}
\label{prop:aks}
\(\Theta\) is Hamiltonian, satisfying \(\{\Theta, \Theta\} = 0\).
\end{proposition}

\begin{proof}
The Poisson bracket is:
\[
\{\Theta, \Theta\} = 2 \iota_{\mathbf{v}} d\Theta = 2 \iota_{\mathbf{v}} d \left( \iota_{\mathbf{v}} \omega_0 + D g^{ab} \partial_a S \partial_b S \right).
\]
Compute:
\[
d \Theta = d \left( \iota_{\mathbf{v}} \omega_0 + D g^{ab} \partial_a S \partial_b S \right).
\]
Since \(\mathcal{L}_{\mathbf{v}} \omega_0 = dS\), and:
\[
d (g^{ab} \partial_a S \partial_b S) = g^{ab} \partial_a (d S \partial_b S) = g^{ab} (d S \partial_b S + \partial_a S d \partial_b S) = 0,
\]
as \(d S \partial_b S\) is symmetric and \(g^{ab}\) is symmetric, the bracket vanishes:
\[
\iota_{\mathbf{v}} d \Theta = \iota_{\mathbf{v}} dS = 0,
\]
since \(dS\) is a 1-form and \(\iota_{\mathbf{v}} dS = \mathbf{v} \cdot dS = 0\) for consistent grading \cite{Alexandrov1997}.
\end{proof}

\subsection{RSVP AKSZ Action}

The action is:
\[
S_{\mathrm{AKSZ}} = \int_{T[1]\Sigma} \left( \Phi dS + \mathbf{v} \cdot \nabla S - D g^{ab} \partial_a S \partial_b S \right).
\]
Integrating over odd coordinates yields:
\[
\mathcal{L}_{\mathrm{RSVP}} = \Phi (\partial_t S + \mathbf{v} \cdot \nabla S) - D |\nabla S|^2.
\]
Variations recover \eqref{eq:PDE1}--\eqref{eq:PDE2}:
\[
\frac{\delta S_{\mathrm{AKSZ}}}{\delta \Phi} = \partial_t S + \mathbf{v} \cdot \nabla S, \quad \frac{\delta S_{\mathrm{AKSZ}}}{\delta S} = -\partial_t \Phi - \nabla \cdot (\Phi \mathbf{v}) + D \nabla^2 S.
\]

\subsection{Path Integral and Boundary Reduction}

Quantization uses the BV path integral:
\[
Z = \int [\mathcal{D}\Psi][\mathcal{D}\Psi^*] \exp\left( \frac{i}{\hbar} S_{\mathrm{BV}} \right).
\]
Fix boundary data \((\Phi, S)|_{\partial \Sigma}\). The boundary current is:
\[
J^\mu = \Phi \partial^\mu S - D \partial^\mu (\Phi S), \quad \delta S_{\mathrm{AKSZ}}|_{\partial \Sigma} = \int_{\partial \Sigma} J^\mu \delta S \, dA_\mu.
\]
The boundary path integral yields:
\[
Z_{\partial \Sigma}[\Phi_{\partial}, S_{\partial}] = \exp\left( \frac{i}{\hbar} \int_{\partial \Sigma} \Phi_{\partial} \Delta S_{\partial} \, dA \right).
\]

\subsection{Discretization and Emergent Transition Kernel}

Partition \(\partial \Sigma\) into cells \(\{\Omega_i\}\). The boundary action is:
\[
S_{\partial} = \sum_{i,j} \Phi_i \left( \Delta t \mathbf{v}_{ij} \cdot \nabla S_j - D \Delta t |\nabla S_j|^2 \right).
\]
The transition amplitude is:
\[
A_{ij} = \exp\left[ \frac{i}{\hbar} \left( \Delta t \mathbf{v}_{ij} \cdot \nabla S_j - D \Delta t |\nabla S_j|^2 \right) \right], \quad U_{ij} = \frac{1}{\sqrt{\mathcal{N}}} A_{ij},
\]
with \(\mathcal{N} = \sum_{ij} |A_{ij}|^2\). The kernel is:
\[
\Gamma_{ij} = |U_{ij}|^2 = \frac{1}{\mathcal{N}} \exp\left[ -\frac{2 D \Delta t}{\hbar} |\nabla S_j|^2 \right].
\]
This is positive and doubly stochastic under normalization.

\begin{proposition}
\label{prop:kernel}
The discretized boundary kernel \(\Gamma_{ij}\) is unistochastic.
\end{proposition}

\begin{proof}
The amplitude \(A_{ij}\) includes the phase \(\theta_{ij} = \Delta t \mathbf{v}_{ij} \cdot \nabla S_j\). Normalizing ensures:
\[
\sum_{ij} |U_{ij}|^2 = \frac{1}{\mathcal{N}} \sum_{ij} |A_{ij}|^2 = 1.
\]
For small \(\Delta t\), the phase term dominates, and \(U_{ij}\) is unitary to \(O(\Delta t^2)\). Thus:
\[
\Gamma_{ij} = |U_{ij}|^2 = \frac{|A_{ij}|^2}{\mathcal{N}},
\]
satisfying the unistochastic condition, as \(\sum_i \Gamma_{ij} = \sum_j \Gamma_{ij} = 1\). The phase \(\theta_{ij}\) ensures consistency with Theorem \ref{thm:unitaryembedding} \cite{Zyczkowski1994}.
\end{proof}

\subsection{Semiclassical Expansion}

Expanding around a classical solution \(\Psi_0 = (\Phi_0, \mathbf{v}_0, S_0)\):
\[
S = S_0 + \hbar^{1/2} \delta S + O(\hbar),
\]
the one-loop correction to the kernel is:
\[
\Gamma_{ij}^{(1)} = \Gamma_{ij}^{(0)} \left[ 1 - \frac{\hbar D \Delta t}{2} (\nabla^2 S_0)^2 + O(\hbar^2) \right].
\]
The fluctuation operator is:
\[
\mathcal{K}_{AB} = \frac{\delta^2 S_{\mathrm{RSVP}}}{\delta \Psi^A \delta \Psi^B}, \quad \mathcal{K}_{\Phi S} = \partial_t S + \mathbf{v} \cdot \nabla S, \quad \mathcal{K}_{S S} = -D \nabla^2 \Phi.
\]
The determinant is regularized via \(\zeta\)-function:
\[
\ln \det \mathcal{K} = -\zeta'(0).
\]

\begin{proof}
Compute second variations:
\[
\frac{\delta^2 S_{\mathrm{RSVP}}}{\delta \Phi \delta S} = \partial_t S + \mathbf{v} \cdot \nabla S, \quad \frac{\delta^2 S_{\mathrm{RSVP}}}{\delta S \delta S} = -D \nabla^2 \Phi.
\]
For the path integral:
\[
Z = \int e^{\frac{i}{\hbar} S_{\mathrm{RSVP}}[\Psi_0 + \delta \Psi]} \approx e^{\frac{i}{\hbar} S_{\mathrm{RSVP}}[\Psi_0]} \int e^{\frac{i}{\hbar} \frac{1}{2} \delta \Psi^A \mathcal{K}_{AB} \delta \Psi^B} = e^{\frac{i}{\hbar} S_{\mathrm{RSVP}}[\Psi_0]} (\det \mathcal{K})^{-1/2}.
\]
Using \(\zeta\)-function regularization, \(\ln \det \mathcal{K} = -\zeta'(0)\), the one-loop correction includes terms proportional to \((\nabla^2 S_0)^2\), derived from the Laplacian’s contribution to \(\mathcal{K}_{S S}\) \cite{Pantev2013}.
\end{proof}

\subsection{Interpretation}

The unistochastic kernel emerges as the boundary manifestation of RSVP’s quantization, with quantum stochasticity as discrete sampling of entropic geometry. The BV–BRST formalism generalizes probability conservation: the master equation \((S_{\mathrm{BV}}, S_{\mathrm{BV}}) = 0\) ensures that stochastic transitions preserve total probability, while entropy production (\(dS\)) measures deviations from exact unitarity, analogous to gauge invariance in field theory \cite{Alexandrov1997}. For non-specialists, this symmetry ensures that probability flows remain consistent across scales, with ghosts and antifields enforcing conservation in the presence of stochastic fluctuations.

\section{Entropy–Curvature Coupling and the Born Rule as a Limit Theorem}

\subsection{Entropy–Curvature Interaction}

On a curved \(\Sigma\) with metric \(g_{\mu\nu}\), scalar curvature \(R\), and Riemann tensor \(R_{\mu\nu\rho\sigma}\), the Lagrangian includes curvature-coupling terms:
\[
\mathcal{L}_{\text{curv}} = \alpha_1 R S^2 + \alpha_2 R_{\mu\nu} \nabla^\mu S \nabla^\nu S + \alpha_3 R |\nabla S|^2,
\]
where \(\alpha_i\) are dimensionless coupling constants. The total Lagrangian is:
\[
\mathcal{L}_{\text{tot}} = \Phi (\partial_t S + \mathbf{v} \cdot \nabla S) - D |\nabla S|^2 + \mathcal{L}_{\text{curv}}.
\]
Curvature regulates coherence: high \(R\) enhances entropy production, accelerating equilibration. The coefficients \(\alpha_i\) arise from one-loop quantum corrections, as computed below.

\subsection{Curvature Correction to the Transition Kernel}

The curvature-corrected kernel is:
\[
\Gamma_{ij}^{(\mathrm{curv})} = |U_{ij}|^2 \exp\left[ -\frac{2 D \Delta t}{\hbar} |\nabla S_j|^2 - \hbar \left( \alpha_1 R S_j^2 + \alpha_3 R |\nabla S_j|^2 \right) \right].
\]
The \(\alpha_i\) terms are derived from one-loop corrections:
\[
\alpha_1 = \frac{\hbar}{24 \pi^2}, \quad \alpha_3 = \frac{\hbar}{12 \pi^2}, \quad \alpha_2 = 0 \text{ (for divergence-free } \mathbf{v}\text{)}.
\]

\begin{proof}
The curvature terms arise from commutators:
\[
[\nabla_\mu, \nabla_\nu] S = R_{\mu\nu\rho\sigma} \partial^\rho S g^{\sigma\tau}.
\]
In the path integral, expand around a classical solution \(\Psi_0\). The one-loop effective action includes:
\[
\Gamma_{\mathrm{anom}} = \frac{\hbar}{24 \pi^2} \int_\Sigma \left( \alpha_1 R \Phi_0 + \alpha_2 T_{\mu\nu\rho} \mathbf{v}_0^\mu \nabla^\rho S_0 + \alpha_3 |\nabla S_0|^2 R \right) d^4x,
\]
where \(T_{\mu\nu\rho}\) is the torsion tensor. For divergence-free \(\mathbf{v}\), \(\alpha_2 = 0\). Compute \(\alpha_1, \alpha_3\) via \(\zeta\)-function regularization:
\[
\ln \det \mathcal{K} = -\zeta'(0).
\]
The fluctuation operator \(\mathcal{K}_{AB}\) includes curvature-dependent terms from \(\mathcal{L}_{\text{curv}}\). Loop integrals over \(\Phi\) and \(S\) propagators yield:
\[
\alpha_1 = \frac{\hbar}{24 \pi^2} \int \frac{d^d k}{(2\pi)^d} \frac{R}{k^2}, \quad \alpha_3 = \frac{\hbar}{12 \pi^2} \int \frac{d^d k}{(2\pi)^d} \frac{|\nabla S|^2 R}{k^2},
\]
regularized in \(d = 4 - \epsilon\) dimensions \cite{Pantev2013}. The exponential form of \(\Gamma_{ij}^{(\mathrm{curv})}\) follows from the boundary action modified by these corrections.
\end{proof}

\subsection{Long-Time Limit and Ergodic Averaging}

The probability evolves:
\[
\rho_i(t + \Delta t) = \sum_j \Gamma_{ij}^{(\mathrm{curv})} \rho_j(t).
\]
For large \(N\):
\[
\rho_i(t + N \Delta t) = \left( \Gamma^{(\mathrm{curv})} \right)^N_{ij} \rho_j(t_0).
\]
By the Perron–Frobenius theorem, since \(\Gamma^{(\mathrm{curv})}\) is positive and irreducible, there exists a unique stationary distribution:
\[
\rho_i^{(\mathrm{eq})} = \sum_j \Gamma_{ij}^{(\mathrm{curv})} \rho_j^{(\mathrm{eq})} = \sum_j |U_{ij}|^2 \rho_j^{(\mathrm{eq})},
\]
as the curvature terms do not affect the leading-order unistochastic structure.

\begin{theorem}[Born Rule as Entropic Limit]
\label{thm:bornlimit}
For an ergodic \(\Gamma^{(\mathrm{curv})}\), empirical frequencies \(f_i^{(N)} = \frac{1}{N} \sum_{k=1}^N \chi_i(t_k)\), where \(\chi_i(t_k) = 1\) if the system is in \(\Omega_i\) at step \(k\), converge almost surely to:
\[
\rho_i^{(\mathrm{eq})} = \sum_j |U_{ij}|^2 \rho_j^{(\mathrm{eq})}.
\]
\end{theorem}

\begin{proof}
Since \(\Gamma^{(\mathrm{curv})}\) is stochastic and ergodic (positive entries ensure irreducibility), the associated Markov chain satisfies Birkhoff’s ergodic theorem:
\[
\lim_{N \to \infty} \frac{1}{N} \sum_{k=1}^N \Gamma^k = \Pi,
\]
where \(\Pi\) is the projection onto the stationary subspace, satisfying \(\Gamma^{(\mathrm{curv})} \rho^{(\mathrm{eq})} = \rho^{(\mathrm{eq})}\). The leading term \(\Gamma_{ij} = |U_{ij}|^2\) ensures the stationary distribution obeys the Born rule:
\[
\rho_i^{(\mathrm{eq})} = \sum_j |U_{ij}|^2 \rho_j^{(\mathrm{eq})}.
\]
The curvature corrections (\(\alpha_1, \alpha_3\)) modify transition rates but preserve stochasticity. The empirical frequencies \(f_i^{(N)}\) converge almost surely to \(\rho_i^{(\mathrm{eq})}\) by the strong law of large numbers for ergodic Markov chains \cite{Risken1996}.
\end{proof}

\subsection{Interpretation}

The Born rule emerges as a statistical attractor of RSVP dynamics, with curvature terms accelerating equilibration. Entropy production, driven by \(\sigma(S) = -\kappa S^2\), measures deviations from microscopic unitarity, while unistochasticity ensures local symplectic conservation, preserving phase coherence within cells \cite{Pantev2013]. This contrasts with Bohmian mechanics, where non-local pilot waves enforce deterministic trajectories, and GRW’s stochastic collapses, which break unitarity.

\section{Cosmological Extension: The Unistochastic Horizon and Expyrotic Reset}

\subsection{From Local Indivisibility to Cosmic Smoothing}

Microscopic indivisible transitions in ISQM (\(\Gamma_{ij} = |U_{ij}|^2\)) aggregate into RSVP’s continuous entropic flow. The cosmic horizon acts as the boundary \(\partial \Sigma\), where local unistochastic transitions translate into global thermodynamic gradients. This scales ISQM’s stochastic jumps to macroscopic smoothing cycles.

\subsection{Entropy Production and Horizon Geometry}

Total entropy is:
\[
S_{\text{tot}}(t) = \int_\Sigma \Phi S \, d^3x, \quad \frac{d S_{\text{tot}}}{dt} = \int_\Sigma \sigma(S) \, d^3x = -\kappa \int S^2 \, d^3x \leq 0.
\]
The entropic Hubble parameter is:
\[
H_{\mathrm{RSVP}} = \frac{1}{3} \nabla \cdot \mathbf{v}.
\]
High \(H_{\mathrm{RSVP}}\) indicates rapid smoothing (analogous to cosmic expansion), while \(H_{\mathrm{RSVP}} \approx 0\) signals entropy traps (e.g., galaxies).

\begin{proof}
From \eqref{eq:PDE1}, integrate:
\[
\frac{d}{dt} \int \Phi S \, d^3x = \int \Phi \left( D \nabla^2 S - \mathbf{v} \cdot \nabla S + \lambda \Phi \right) + S \left( -\kappa S^2 - \nabla \cdot (\Phi \mathbf{v}) \right) d^3x.
\]
After integration by parts and applying boundary conditions, the dominant term is \(-\kappa \int S^3 \, d^3x \leq 0\), ensuring non-positive entropy change \cite{Seifert2012].
\end{proof}

\subsection{The Unistochastic Horizon}

At cosmological scales, coherence degrades across the particle horizon. The horizon kernel is:
\[
\Gamma_{ij}^{(\mathrm{cos})} = |U_{ij}|^2 \exp\left[ -\frac{\Lambda_{\mathrm{ent}} d_{ij}^2}{\hbar c} \right],
\]
where \(d_{ij}\) is the comoving separation, and \(\Lambda_{\mathrm{ent}}\) is an entropic cosmological constant, analogous to the cosmological constant in general relativity.

\subsection{Entropy–Curvature Feedback and the Expyrotic Limit}

Curvature-entropy coupling reduces gradients:
\[
\nabla S \to 0, \quad R \to 0, \quad H_{\mathrm{RSVP}} \to 0,
\]
approaching an Expyrotic state of maximal smoothness, where entropy production ceases.

\subsection{Quantum Seeds and the Reset Mechanism}

At Expyrosis, the kernel becomes uniform:
\[
\Gamma_{ij} \to \frac{1}{N}.
\]
Fluctuations \(\epsilon_{ij}\) in \(\Gamma_{ij} = \frac{1}{N} (1 + \epsilon_{ij})\) grow as:
\[
\epsilon_{ij}(t) \sim e^{\lambda t}, \quad \lambda = \left. \frac{\partial \sigma}{\partial S} \right|_{S=0} = 2 \kappa S.
\]

\begin{proof}
Linearize around \(\Gamma_{ij} = \frac{1}{N}\):
\[
\dot{\epsilon}_{ij} = \frac{\partial \sigma}{\partial S} \epsilon_{ij} = 2 \kappa S \epsilon_{ij}.
\]
For small \(S\), the exponential growth seeds new structure, as fluctuations break the uniform state, initiating a new entropic cycle \cite{Seifert2012}.
\end{proof}

\subsection{Cosmological Born Rule}

Macrostate probabilities follow:
\[
P(\mathcal{O}_i) = \lim_{T \to \infty} \frac{1}{T} \int_0^T |U_{ij}(t)|^2 \, dt.
\]
This generalizes the Born rule to cosmological scales, predicting distributions for CMB anisotropies and galaxy clustering.

\subsection{Interpretation}

The universe evolves through entropic cycles, with stochastic resets seeding new structure. The unistochastic horizon replaces spatial expansion with probabilistic attenuation, unifying quantum and cosmological dynamics.

\subsection{Summary}

\begin{itemize}[noitemsep]
\item The cosmic horizon is the boundary of unistochastic coherence.
\item Curvature-entropy feedback drives Expyrosis.
\item Fluctuations in \(\Gamma_{ij}\) trigger stochastic resets.
\item The Born rule extends to cosmological ensembles.
\end{itemize}

\section{Unified View: Indivisibility Across Scales}

\subsection{Indivisibility as a Universal Principle}

Indivisibility resists arbitrary partition. In ISQM, it is the non-factorizable kernel \(\Gamma_{ij}(t | h)\), preventing decomposition into intermediate states. In RSVP, it is the continuous entropy flux, ensuring no information is lost. Cosmologically, it is the coherent plenum, where boundaries are temporary entropy discontinuities.

\subsection{Scale Hierarchy of the Plenum}

\begin{center}
\begin{tabular}{@{}lll@{}}
\toprule
\textbf{Scale} & \textbf{Dominant Structure} & \textbf{Conservation Law} \\
\midrule
Quantum & Unistochastic kernel \(|U_{ij}|^2\) & Unitarity (phase coherence) \\
Mesoscopic & RSVP field \((\Phi, \mathbf{v}, S)\) & Entropy continuity \\
Macroscopic & Entropy curvature \(R, S\) & Thermodynamic balance \\
Cosmic & Horizon kernel \(\Gamma_{ij}^{(\mathrm{cos})}\) & Global entropy conservation \\
\bottomrule
\end{tabular}
\end{center}

\subsection{The RSVP–UQM Correspondence Principle}

The functor \(\mathfrak{C}_1: \mathbf{RSVPField} \to \mathbf{UniStoch}\) and adjoint \(\mathfrak{R}_1\) define:
\[
\mathfrak{C}_1(\Phi, \mathbf{v}, S) = |U|^2, \quad \mathfrak{R}_1(|U|^2) = (\Phi, \mathbf{v}, S).
\]
This unifies ISQM’s discrete stochasticity with RSVP’s continuous geometry.

\subsection{From Quantization to Smoothing}

The deformation parameters are:
\[
\hbar \text{ (quantum)}, \quad D \text{ (smoothing)}, \quad \Lambda_{\mathrm{RSVP}} = \frac{D}{\hbar}.
\]
The BV–BRST formalism ensures probability conservation, with entropy production as a deformation.

\subsection{Expyrotic Cycles as Indivisible Renewal}

Expyrotic resets re-quantize the universe via stochastic fluctuations, initiating new entropic cycles.

\subsection{Philosophical Implications}

\begin{enumerate}[noitemsep]
\item \textbf{Reality as recursive coherence}: Existence is the persistence of entropy-preserving patterns.
\item \textbf{Measurement as entropic coupling}: Observation equilibrates coherence via entropy exchange.
\item \textbf{Time as smoothing gradient}: Temporal order emerges from entropy equalization.
\item \textbf{Cosmos as indivisible field}: Boundaries are temporary entropy discontinuities.
\end{enumerate}

\subsection{Outlook}

\begin{itemize}[noitemsep]
\item Derived Markov stacks for higher-categorical formulations.
\item Numerical RSVP simulations to verify Born rule convergence.
\item Cosmological data fitting with CMB anisotropies.
\end{itemize}

\subsection{Closing Reflection}

The universe equilibrates, not expands. The wave smooths, not collapses. Reality reconstitutes, not divides.

\begin{center}
\emph{The universe does not expand; it equilibrates. The wave does not collapse; it smooths. Reality does not divide; it reconstitutes.}
\end{center}

\section{Double-Slit Experiment in ISQM and RSVP}

\subsection{Setup}

Consider a double-slit experiment with slits at positions \(x_1, x_2\) in a 1D configuration space \(X\), and a detector screen at positions \(\{y_i\}\). In ISQM, the system transitions from an initial state to slit positions \(x_1, x_2\), then to detector positions \(y_i\), with probabilities \(\Gamma_{ij}(t | h)\) dependent on the slit history \(h\). In RSVP, the slits are modeled as cells \(\Omega_1, \Omega_2\) in \(\Sigma\), with \(\Phi\) peaked at \(x_1, x_2\), and the detector as cells \(\{\Omega_i\}\).

\subsection{ISQM Prediction}

The transition matrix from slits to detectors is:
\[
\Gamma_{ij} = |U_{ij}|^2, \quad U = \frac{1}{\sqrt{2}} \begin{pmatrix} e^{i \theta_1} & e^{i \theta_2} \\ e^{i \theta_2} & e^{i \theta_1} \end{pmatrix},
\]
where \(\theta_1, \theta_2\) are phases associated with paths through slits \(x_1, x_2\). The probability at detector \(y_i\) is:
\[
\rho_i = \sum_{j=1}^2 \Gamma_{ij} \rho_j = \frac{1}{2} \left( |e^{i \theta_1}|^2 + |e^{i \theta_2}|^2 + 2 \Re (e^{i (\theta_1 - \theta_2)}) \right) \rho_j,
\]
assuming equal initial probabilities \(\rho_1 = \rho_2 = 1/2\). The interference term \(2 \Re (e^{i (\theta_1 - \theta_2)})\) produces the characteristic double-slit pattern \cite{Barandes2023}.

\subsection{RSVP Prediction}

Model \(\Sigma\) as a 1D spatial domain with cells \(\Omega_1, \Omega_2\) at \(x_1, x_2\), and detector cells \(\Omega_i\). Set \(\Phi(x) = \frac{1}{2} (\delta(x - x_1) + \delta(x - x_2))\), \(\mathbf{v} = v_0 \hat{x}\), \(S = -\ln \Phi + S_0\), with \(v_0 = 1\), \(D = 0.01\), \(\Delta t = 0.1\). The flux across cell boundaries is:
\[
J_{ij} = \int_{\partial \Omega_i \cap \Omega_j} (\Phi v_0 - D \nabla \Phi) \, dx.
\]
For simplicity, assume \(D \approx 0\), so:
\[
\Gamma_{ij} = \frac{1}{\Phi_j} \int_{\partial \Omega_i \cap \Omega_j} \Phi v_0 \, dx.
\]
Discretize the detector into \(N\) cells, with \(\Phi_j = 1/N\). The transition matrix is:
\[
\Gamma_{ij} = \frac{v_0 \Delta t}{\Phi_j}, \quad \sum_i \Gamma_{ij} = 1.
\]
Set \(\mathbf{v} = \nabla \theta\), with \(\theta(x) = k (x - x_j)\), yielding phases \(\theta_{ij} = k (y_i - x_j)\). The unitary is:
\[
U_{ij} = \sqrt{\Gamma_{ij}} e^{i k (y_i - x_j) / \hbar}.
\]
The probability distribution:
\[
\rho_i = \sum_j |U_{ij}|^2 \rho_j = \frac{1}{2} \sum_j \Gamma_{ij} \left( 1 + \cos \left( \frac{k (y_i - x_1) - k (y_i - x_2)}{\hbar} \right) \right),
\]
reproducing the interference pattern, consistent with ISQM and standard quantum mechanics.

\subsection{Numerical Example}

Discretize the screen into 100 cells, with \(y_i = i \Delta y\), \(\Delta y = 0.01\), \(x_1 = -0.5\), \(x_2 = 0.5\), \(k = 2\pi\), \(\hbar = 1\). Compute:
\[
\Gamma_{ij} = \frac{0.1}{1/100} = 0.01, \quad U_{ij} = \sqrt{0.01} e^{i 2\pi (i \cdot 0.01 - x_j)}.
\]
Sum probabilities:
\[
\rho_i \propto \cos^2 \left( \pi (i \cdot 0.01 - 0.5) \right) + \cos^2 \left( \pi (i \cdot 0.01 + 0.5) \right) + 2 \cos \left( \pi (i \cdot 0.01 - 0.5) \right) \cos \left( \pi (i \cdot 0.01 + 0.5) \right).
\]
This yields the double-slit interference pattern, verifiable numerically.

\subsection{Interpretation}

Both ISQM and RSVP reproduce the quantum interference pattern without wavefunctions. ISQM uses history-dependent transitions, while RSVP derives \(\Gamma_{ij}\) from entropic fluxes, unifying discrete and continuous descriptions.

\section{EPR Experiment in ISQM and RSVP}

\subsection{Setup}

Consider an EPR experiment with two entangled particles A and B, each with spin-1/2, measured along directions \(\mathbf{a}\) and \(\mathbf{b}\). The configuration space is \(X = \{ (\uparrow_A, \uparrow_B), (\uparrow_A, \downarrow_B), (\downarrow_A, \uparrow_B), (\downarrow_A, \downarrow_B) \}\). In ISQM, the transition matrix \(\Gamma_{ij}(t | h)\) depends on the entangled history. In RSVP, the system is modeled as cells \(\{\Omega_i\}\) corresponding to spin configurations, with \(\Phi\) representing probability densities.

\subsection{ISQM Prediction}

For a singlet state, the unitary matrix is:
\[
U = \frac{1}{\sqrt{2}} \begin{pmatrix} 1 & 0 & 0 & 1 \\ 0 & 1 & -1 & 0 \\ 0 & -1 & 1 & 0 \\ -1 & 0 & 0 & 1 \end{pmatrix},
\]
yielding:
\[
\Gamma_{ij} = |U_{ij}|^2 = \frac{1}{2} \begin{pmatrix} 1 & 0 & 0 & 1 \\ 0 & 1 & 1 & 0 \\ 0 & 1 & 1 & 0 \\ 1 & 0 & 0 & 1 \end{pmatrix}.
\]
The joint probability for measurements along \(\mathbf{a}, \mathbf{b}\) is:
\[
P(\uparrow_A, \downarrow_B | \mathbf{a}, \mathbf{b}) = \frac{1}{2} \sin^2 \left( \frac{\theta}{2} \right),
\]
where \(\theta\) is the angle between \(\mathbf{a}\) and \(\mathbf{b}\), reproducing quantum correlations \cite{Barandes2023}.

\begin{theorem}[Locality of ISQM Correlations]
\label{thm:locality}
ISQM correlations are local, reproducing Bell inequality violations without superluminal signaling.
\end{theorem}

\begin{proof}
Consider two subsystems A and B with joint configuration \((x_A, x_B) \in X_A \times X_B\). The transition matrix factorizes for independent measurements:
\[
\Gamma_{(i k)(j l)} = \Gamma_{ij}^A \Gamma_{kl}^B,
\]
unless correlated by shared history \(h\). The history \(h\) is determined by local interactions (e.g., entanglement preparation) at light-like or time-like separations. For the singlet state, the conditional probability:
\[
P(x_A | x_B, \mathbf{a}, \mathbf{b}, h) = \frac{\Gamma_{(i k)(j l)} \rho_{j l}}{\sum_m \Gamma_{(i m)(j l)} \rho_{j l}},
\]
depends on \(h\), which encodes the entangled state locally. Bell violations arise from non-Markovianity:
\[
\langle \mathbf{a} \cdot \mathbf{S}_A \mathbf{b} \cdot \mathbf{S}_B \rangle = -\frac{1}{4} \mathbf{a} \cdot \mathbf{b},
\]
matching quantum mechanics without superluminal signaling, as probabilities are fixed by prior local interactions. The no-signaling theorem ensures locality \cite{Barandes2023}.
\end{proof}

\subsection{RSVP Prediction}

Model the system as four cells \(\Omega_1, \ldots, \Omega_4\) corresponding to \((\uparrow_A, \uparrow_B), \ldots, (\downarrow_A, \downarrow_B)\). Set \(\Phi_i = 1/4\), \(\mathbf{v} = \nabla \theta\), with \(\theta_{ij} = \hbar \arg(U_{ij})\). The flux is:
\[
J_{ij} = \frac{1}{4} \int_{\partial \Omega_i \cap \Omega_j} \nabla \theta \cdot d\mathbf{A}.
\]
For simplicity, assume \(D = 0\), and set:
\[
\Gamma_{ij} = \frac{J_{ij}}{\Phi_j} = |U_{ij}|^2.
\]
Using the singlet unitary, compute:
\[
\rho_i = \sum_j \Gamma_{ij} \rho_j = \frac{1}{2} \left( \delta_{i,1} + \delta_{i,4} \right),
\]
for initial \(\rho_j = 1/4\), reproducing the entangled correlations. The phase \(\theta_{ij}\) ensures:
\[
P(\uparrow_A, \downarrow_B) = \frac{1}{2} \sin^2 \left( \frac{\theta}{2} \right),
\]
matching ISQM and quantum mechanics.

\subsection{Numerical Example}

Discretize into four cells, with \(\Phi_i = 0.25\), \(\Delta t = 0.1\), \(\mathbf{v} = (1, 0, 0)\). Compute fluxes numerically, yielding \(\Gamma_{ij}\) as above. Verify:
\[
\sum_i \Gamma_{ij} = \frac{1}{2} (1 + 1) = 1, \quad \sum_j \Gamma_{ij} = 1.
\]
Construct \(U_{ij} = \sqrt{\Gamma_{ij}} e^{i \theta_{ij} / \hbar}\), with \(\theta_{ij}\) from the singlet unitary, confirming unistochasticity.

\subsection{Interpretation}

ISQM and RSVP reproduce EPR correlations without non-locality. ISQM uses history-dependent transitions, while RSVP derives correlations from entropic fluxes, with \(\mathbf{v} = \nabla \theta\) encoding entanglement phases. This contrasts with Bohmian mechanics’ non-local guidance.

\section{Detailed Balance and Reversibility}

\subsection{Definition of Reversible Flows}

A flow is reversible if \(\mathbf{v}(x, t) = -\mathbf{v}(-
